% GENERATED FILE. Edit canonical lessons, then run npm run generate:books.
% canonical-source-hash: fnv1a64:74a5eb03be049269
% canonical-lessons: KA-C218-siddha, KA-C218-labhya, KA-C218-spasta, KA-C218-pramanika, KA-C218-udara

\chapter{Ready, Available, Clear, Honest, Generous}
\label{ch:ka-ready-available-clear-honest-generous}
\hlchaptermodality{218}

\begin{chapteropening}
\textbf{By the end of this chapter:} \emph{I can say ready, available, clear, honest and generous.}

Complete the last lesson of chapter 218: I can say ready, available, clear, honest and generous.
\end{chapteropening}

\section[\texorpdfstring{\kn{ಸಿದ್ಧ} (siddha)}{ಸಿದ್ಧ (siddha)}]{\kn{ಸಿದ್ಧ} (siddha) — ready}

\label{lesson:KA-C218-siddha}



\begin{quote}
Before the new one: say the Kannada for loose, then the Kannada for thick.
\end{quote}

\subsection*{You'll want to know: \kn{ಸಿದ್ಧ}}
\textbf{\kn{ಸಿದ್ಧ}} — \emph{siddha} — \textquotedblleft{}ready\textquotedblright{}.

\begin{grammarlens}[title={What you've built}]
One more word for this chapter.
\end{grammarlens}

\subsection*{Your turn}
\begin{itemize}
\raggedright
  \item \emph{Say it:} \emph{siddha}
  \item \emph{Say it:} \emph{siddha}, once more
  \item \emph{Say it:} say \emph{siddha}
  \item \emph{Recall:} say the Kannada for loose, then the Kannada for thick, then say \emph{siddha} again
\end{itemize}

\begin{tcolorbox}[breakable,colback=teal!4,colframe=teal!35!black,title={Before you move on}]
What does \kn{ಸಿದ್ಧ} mean? (\textquotedblleft{}Ready\textquotedblright{}.) Say it once more.
\end{tcolorbox}
\section[\texorpdfstring{\kn{ಲಭ್ಯ} (labhya)}{ಲಭ್ಯ (labhya)}]{\kn{ಲಭ್ಯ} (labhya) — available}

\label{lesson:KA-C218-labhya}



\begin{quote}
Before the new one: say the Kannada for thick, then the Kannada for ready.
\end{quote}

\subsection*{You'll want to know: \kn{ಲಭ್ಯ}}
\textbf{\kn{ಲಭ್ಯ}} — \emph{labhya} — \textquotedblleft{}available\textquotedblright{}.

\begin{grammarlens}[title={What you've built}]
One more word for this chapter.
\end{grammarlens}

\subsection*{Your turn}
\begin{itemize}
\raggedright
  \item \emph{Say it:} \emph{labhya}
  \item \emph{Say it:} \emph{labhya}, once more
  \item \emph{Say it:} say \emph{labhya}
  \item \emph{Recall:} say the Kannada for thick, then the Kannada for ready, then say \emph{labhya} again
\end{itemize}

\begin{tcolorbox}[breakable,colback=teal!4,colframe=teal!35!black,title={Before you move on}]
What does \kn{ಲಭ್ಯ} mean? (\textquotedblleft{}Available\textquotedblright{}.) Say it once more.
\end{tcolorbox}
\section[\texorpdfstring{\kn{ಸ್ಪಷ್ಟ} (spaṣṭa)}{ಸ್ಪಷ್ಟ (spaṣṭa)}]{\kn{ಸ್ಪಷ್ಟ} (spaṣṭa) — clear}

\label{lesson:KA-C218-spasta}



\begin{quote}
Before the new one: say the Kannada for ready, then the Kannada for available.
\end{quote}

\subsection*{You'll want to know: \kn{ಸ್ಪಷ್ಟ}}
\textbf{\kn{ಸ್ಪಷ್ಟ}} — \emph{spaṣṭa} — \textquotedblleft{}clear\textquotedblright{}.

\begin{grammarlens}[title={What you've built}]
One more word for this chapter.
\end{grammarlens}

\subsection*{Your turn}
\begin{itemize}
\raggedright
  \item \emph{Say it:} \emph{spaṣṭa}
  \item \emph{Say it:} \emph{spaṣṭa}, once more
  \item \emph{Say it:} say \emph{spaṣṭa}
  \item \emph{Recall:} say the Kannada for ready, then the Kannada for available, then say \emph{spaṣṭa} again
\end{itemize}

\begin{tcolorbox}[breakable,colback=teal!4,colframe=teal!35!black,title={Before you move on}]
What does \kn{ಸ್ಪಷ್ಟ} mean? (\textquotedblleft{}Clear\textquotedblright{}.) Say it once more.
\end{tcolorbox}
\section[\texorpdfstring{\kn{ಪ್ರಾಮಾಣಿಕ} (prāmāṇika)}{ಪ್ರಾಮಾಣಿಕ (prāmāṇika)}]{\kn{ಪ್ರಾಮಾಣಿಕ} (prāmāṇika) — honest}

\label{lesson:KA-C218-pramanika}



\begin{quote}
Before the new one: say the Kannada for available, then the Kannada for clear.
\end{quote}

\subsection*{You'll want to know: \kn{ಪ್ರಾಮಾಣಿಕ}}
\textbf{\kn{ಪ್ರಾಮಾಣಿಕ}} — \emph{prāmāṇika} — \textquotedblleft{}honest\textquotedblright{}.

\begin{grammarlens}[title={What you've built}]
One more word for this chapter.
\end{grammarlens}

\subsection*{Your turn}
\begin{itemize}
\raggedright
  \item \emph{Say it:} \emph{prāmāṇika}
  \item \emph{Say it:} \emph{prāmāṇika}, once more
  \item \emph{Say it:} say \emph{prāmāṇika}
  \item \emph{Recall:} say the Kannada for available, then the Kannada for clear, then say \emph{prāmāṇika} again
\end{itemize}

\begin{tcolorbox}[breakable,colback=teal!4,colframe=teal!35!black,title={Before you move on}]
What does \kn{ಪ್ರಾಮಾಣಿಕ} mean? (\textquotedblleft{}Honest\textquotedblright{}.) Say it once more.
\end{tcolorbox}
\section[\texorpdfstring{\kn{ಉದಾರ} (udāra)}{ಉದಾರ (udāra)}]{\kn{ಉದಾರ} (udāra) — generous}

\label{lesson:KA-C218-udara}



\begin{quote}
Before the new one: say the Kannada for clear, then the Kannada for honest.
\end{quote}

\subsection*{You'll want to know: \kn{ಉದಾರ}}
\textbf{\kn{ಉದಾರ}} — \emph{udāra} — \textquotedblleft{}generous\textquotedblright{}.

\begin{grammarlens}[title={What you've built}]
One more word for this chapter.
\end{grammarlens}

\subsection*{Your turn}
\begin{itemize}
\raggedright
  \item \emph{Say it:} \emph{udāra}
  \item \emph{Say it:} \emph{udāra}, once more
  \item \emph{Say it:} say \emph{udāra}
  \item \emph{Recall:} say the Kannada for clear, then the Kannada for honest, then say \emph{udāra} again
\end{itemize}

\begin{tcolorbox}[breakable,colback=teal!4,colframe=teal!35!black,title={Before you move on}]
What does \kn{ಉದಾರ} mean? (\textquotedblleft{}Generous\textquotedblright{}.) Say it once more.
\end{tcolorbox}
